\section{Limitations}
\label{sec:limitations}

Several limitations are known independently of the confirmatory result and are
stated here rather than deferred.

\paragraph{Effective token count measures brevity, not vagueness.}
The informativeness rule that separates informative from low-information context
counts tokens carrying alphanumeric content after de-identification placeholders
are removed. It therefore detects emptiness and brevity. A genuinely vague but
lengthy indication --- ``evaluate'' followed by boilerplate --- scores as
informative. This is a measurement limitation of the state assignment, not of
the controlled interventions, which do not depend on it beyond eligibility.

\paragraph{The primary evaluation partition is not an official split.}
The official test split of the source dataset yields too few patients within the
study cohort to support the primary comparison, so a patient-disjoint partition
was carved from the training split using patient identifiers and a frozen seed
alone. Primary estimates are consequently not directly comparable to published
official-split baselines. The official validation and test splits are retained
and reported as secondary confirmation.

\paragraph{The transport result is not robust to label source.}
H1 is established under the primary label endpoint and reverses under the
pre-specified sensitivity endpoint. The two findings-derived label sets proved
non-commensurable across sites, differing by a factor of five in labelling
density, so the sensitivity cannot arbitrate the cross-site contrast in either
direction. What this leaves is a transport failure demonstrated for one label
endpoint, with its label-source robustness untested rather than refuted.
Establishing it would require a label endpoint that is dense and comparable at
both institutions, which neither available endpoint is.

\paragraph{Two training seeds, and only for two of the four models.}
The bootstrap quantifies uncertainty in \emph{evaluation}, not variability
across training runs. A second seed was run for the image-only and
feature-fusion models and reproduced every verdict
(\cref{sec:seed-replicate}), but two runs bound seed variability loosely: point
estimates moved by up to a third of their own magnitude while keeping their
sign and their verdict. Differences between models smaller than the materiality
threshold remain unestablished against seed-to-seed variation. The text-only
model was not retrained, so neither it nor the late-fusion model that consumes
it has been replicated.

\paragraph{Source images are lossy, external images are lossless.}
The source dataset publishes JPEG images and the external dataset publishes PNG.
Both were reduced to the model input resolution by a single identical bilinear
step, so no additional resampling asymmetry was introduced, but the compression
history of the two sites differs at origin. This difference tracks the site
variable and cannot be removed by preprocessing choices.

\paragraph{One external image is unusable at source.}
One frontal image in the external dataset matches the publisher's recorded
checksum yet fails to decode: its PNG header parses while an IDAT chunk fails
its cyclic redundancy check. The damage is upstream and no retransfer repairs
it. That study is excluded, and the exclusion is reported rather than absorbed.

\paragraph{Studies whose only images lack a view label are excluded.}
Images carrying no view label are treated as non-frontal, which is the
conservative reading, since assuming them frontal would admit views the protocol
excludes. Such exclusions may not be missing at random.

\paragraph{Section boundaries derive from a curated header vocabulary.}
Reports whose headers fall outside that vocabulary are excluded rather than
parsed heuristically, which trades coverage for the guarantee that no
post-diagnostic text enters a pre-diagnostic field.

\paragraph{No sociodemographic subgroup or fairness analysis was performed.}
This study reports no patient demographics and no subgroup performance
estimates, and applied no fairness-specific method. The omission was not a
considered exclusion: the protocol specified cohorts, endpoints, contrasts and
materiality without naming sociodemographic strata, so no such analysis was
frozen in advance, and adding one after the external results were inspected
would be post-hoc exploratory by the protocol's own rule.

The gap matters here more than it would in a study reporting only aggregate
accuracy, because it interacts with the central finding. If a frozen abstention
policy transports poorly between institutions, there is no reason to assume it
transports uniformly across patient groups within one. A policy that deferred
disproportionately on one group would be a fairness failure invisible to every
metric reported here --- the same silence that makes coverage monitoring
inadequate for institutional transfer would hide it. Both datasets publish
demographic fields sufficient to test this, and doing so is the natural
companion to the third site identified in \cref{sec:future-work}.
